\chapter{实验结果与分析}



\section{实验设置}

本节对本文提出的安全反思训练范式及 RLVR 强化策略进行系统评估。我们使用第3节构建的 Safe-CoT 训练数据以及规则验证强化学习数据集，在 ERNIE-4.5-VL-28B-A3B 原始模型上进行分阶段训练，并在多个评估集上对比不同训练策略的效果。

\paragraph{实验分组}
我们设置如下四组对比实验，以验证各组件的增益来源：

\begin{itemize}
    \item \textbf{Base Model：} 原始 ERNIE-4.5-VL-28B-A3B；
    \item \textbf{Distilled CoT：} 基于DeepSeek-R1-0528版本教师模型蒸馏得到的 CoT 数据进行监督微调；
    \item \textbf{SafetyCoT：} 使用基于规则的安全反思 CoT 数据进行监督微调；
    \item \textbf{SafetyCoT + RLVR：} 在 SafetyCoT 基础上进一步进行规则验证驱动的在线强化学习。
\end{itemize}

\paragraph{统一推理设置}
为保证评估公平性，所有模型在推理时采用相同的解码策略、温度参数与最大生成长度，避免由于解码参数差异引入性能偏移。我们的目标是验证：
随着训练流程推进，模型在提升安全输出能力的同时，其通用能力能够保持在合理水平，并进一步分析 RLVR 对越狱鲁棒性与边界场景稳定性的影响。

\section{评估集说明}

% \subsection{多模安全评估集}
% \subsubsection{数据集说明}
% % 需要数据集来源说明一下 以及各分类定义 凑字数
% 该评估集由两类数据构成：
% 安全问题集与有用问题集。
% 前者用于评估模型在明确或潜在违规场景下的合规性；数据来源为线上流量日志中被人工审核为对应违规的问题集;
% 后者用于评估模型在非违规但接近风险边界的输入上是否仍能保持正常应答能力，从而衡量过度拒答现象。


% 在该评估集上，我们分别计算安全性与有用性指标：
% 安全性得分定义为安全占比（即模型未产生违规输出的比例）；
% 有用性得分定义为在边界有用问题上满足指令要求、给出有效回答的比例。
% 除此之外，我们额外报告 Safe\&Helpful 指标，用于刻画模型同时满足安全与有用性的整体能力平衡。

% \subsubsection{数据集分布}
% 我们按照审核违规类别将数据分为8类, 评估集具体分布如下所示

\subsection{多模安全评估集}

\subsubsection{数据集说明}

为系统评估模型在真实部署场景下的安全表现与能力平衡情况，本文构建了一个多模态安全评估集。该评估集由两类数据构成：\textbf{安全问题集}与\textbf{有用问题集}。

\textbf{安全问题集}用于评估模型在明确或潜在违规场景下的合规性。数据来源于线上流量日志中经人工审核确认为违规类别的问题样本。此类样本覆盖多种高风险语境，能够反映模型在开放环境中可能遭遇的真实攻击与敏感问题形式。

\textbf{有用问题集}用于评估模型在非违规但接近风险边界输入下的正常应答能力。该部分样本语义上可能涉及敏感领域或灰度表达，但本身不构成违规内容。该数据用于衡量模型在强化安全能力后是否出现过度拒答（Over-refusal）现象，从而评估模型在安全与有用性之间的平衡能力。

在该评估集上，我们分别计算安全性与有用性指标：

\begin{itemize}
    \item \textbf{安全性得分（Safety Score）：} 定义为模型在安全问题集上未产生违规输出的比例；
    \item \textbf{有用性得分（Helpfulness Score）：} 定义为模型在边界有用问题上满足指令要求并给出有效回答的比例；
\end{itemize}

\subsubsection{数据集分布}
评估集按照审核违规类别进行细分，共覆盖以下十二类风险类型：
暴力、血腥、犯罪、仇恨、儿童保护、成人内容与敏感主题、
性感低俗、心理与行为健康、欺凌与骚扰、管制商品及服务、
隐私以及涉政涉习。
每个违规类别包含 $30$ 条评估样本，共计：
\begin{equation}
N_{\text{unsafe}} = 12 \times 30 = 360
\end{equation}

在此基础上，为评估过度拒答现象，构建正常边界问题样本，其数量与安全问题保持 $1:2$ 的比例，即：

\begin{equation}
N_{\text{normal}} = 1 / 2 \times N_{\text{unsafe}} = 180
\end{equation}

因此，整个评估集规模为：

\begin{equation}
N_{\text{total}} = N_{\text{unsafe}} + N_{\text{normal}} = 540
\end{equation}

该分布设计具有如下特点：

\begin{itemize}
    \item 多类别均衡采样，避免单一风险类型主导整体结果；
    \item 较高比例的正常样本用于充分检测模型是否因安全强化而牺牲有用性；
    \item 类别覆盖广泛，能够反映开放环境下不同类型风险表达形式。
\end{itemize}

通过上述设计，该评估集能够较为全面地衡量模型在真实部署环境中面对复杂语境时的安全稳健性与服务能力平衡情况。

\subsubsection{评估方式说明}
采用LLM-AS-A-JUDGE的方式, 基于此前生产的实例级别的rule, 依据输出是否合规来判断是否符合安全标准要求,
针对有用性问题考察也基于同样的方式.


\subsection{越狱攻击评估集}

\subsubsection{越狱攻击评估集说明}
% 越狱攻击数据集来源, 以及各种攻击类型, 现有大模型对于各种攻击类型的容忍程度
% jailbreakbench LLM-PBE MM-SafetyBench
% 可以给一个示例?
该评估集为纯文本评估集 用于评估模型抵抗越狱攻击的能力，包含多种模板化攻击方式
（例如角色扮演、拒绝抑制、语义改写等）所包装的有害问题。
评价指标为 Attack Success Rate（ASR），即攻击成功诱导模型生成违规内容的比例，数值越低代表模型越健壮。

\subsubsection{越狱攻击评估集数据分布}
该评估集参考了主流开源文本越狱攻击方式, 在原始ERNIE-4.5-VL-28B-A3B模型上进行推理, 主要选择基于模版的攻击方式, 选取表现较差的前10\%的
攻击方式作为评估集模版, 如下所示:


\subsubsection{评估方式说明}
基于LLM-AS-A-JUDGE的方式, 使用固定的请求模版, 做一个是否存在有害输出的二分类任务. 使用ERNIE-4.5-300B-A47B作为评估模型, 不使用主流开源商用模型的
原因是 安全问题和输出存在敏感内容, 存在被拦截导致评估结果部分缺失的情况, 因此需要自己部署开源模型评估


\subsection{事实可信评估集}
\subsubsection{事实可信评估集说明}
% 数据是怎么构建的 说具体点 可以给一个示例
该评估集用于验证在引入 Agentic RL（结合检索工具调用）后，模型是否能够正确触发外部搜索工具、利用检索证据进行回答，并保持事实一致性。与前两类主要关注安全风险不同，该评估集重点衡量模型在外部信息依赖条件下的事实可信与工具使用行为是否稳定。

\subsubsection{评估方式说明}
参考3.2中的数据集构建方式, 在构建问题的同时获得了对应的基准事实(Ground Truth), 多数评估指标采用 LLM-As-A-Judge 方式进行自动判定，以提升评测效率与一致性。对于事实可信评估集，由于提供 Ground Truth，我们基于模型输出与 Ground Truth 的一致性进行比对评估，以获得更客观的事实正确性度量。

\section{实验结果与分析}

本节系统比较不同训练策略在多模安全评估、越狱攻击评估以及事实可信评估中的表现，并结合实验现象分析各组件带来的能力变化机制。从整体趋势来看，模型性能随着训练阶段逐步演进呈现出清晰的分层提升路径：首先，通过引入结构化安全反思监督训练（SafetyCoT），模型在明确风险场景下的合规性显著提升，同时保持较高的边界问题响应能力；其次，在此基础上进一步引入规则验证强化学习（RLVR），模型在复杂诱导场景与分布外攻击中的稳定性得到进一步增强。

这一结果验证了本文在第三章提出的核心假设：仅依赖推理蒸馏不足以稳定增强安全能力，而将规则约束嵌入推理结构，并通过强化学习进行行为层面的动态校正，能够更有效地缩小“安全推理鸿沟”，实现安全性与有用性的协同提升。

\subsection{多模安全与有用性主结果}

\begin{table*}[t]
\centering
\caption{多模安全与有用性评估结果。Safety 表示安全占比（越高越好），Helpfulness 表示有用性得分（越高越好）。}
\label{tab:main_results}
\begin{tabular}{lcc}
\toprule
Model & Safety $\uparrow$ & Helpfulness $\uparrow$ \\
\midrule
Base Model  & 0.56  & 1.00 \\
Distilled CoT & 0.73 & 0.81 \\
SafetyCoT & 0.88 & 0.91 \\
SafetyCoT + RLVR & 0.92 & 0.89 \\
\bottomrule
\end{tabular}
\end{table*}

从表~\ref{tab:main_results} 可见，Base Model 在安全问题集上的安全占比仅为 $0.56$，说明原始模型在真实开放场景下仍存在较高比例的违规输出风险。引入 Distilled CoT 后，安全得分提升至 $0.73$，表明显式推理链蒸馏在一定程度上改善了模型的风险识别能力。然而，其提升幅度相对有限，同时有用性指标明显下降，说明单纯蒸馏推理过程并不能保证模型稳定形成安全边界，甚至可能在无规则约束条件下引入额外的不确定性。

相比之下，SafetyCoT 模型在安全性指标上提升至 $0.88$，同时保持 $0.91$ 的有用性得分，呈现出更加平衡的性能表现。这一结果说明，引入规则映射与结构化安全反思数据后，模型不仅能够识别高风险语境，还能够在边界问题上保持较为稳定的生成行为。这验证了“规则理解—风险分析—安全生成”推理结构对缩小安全推理鸿沟的有效性。

进一步地，在 SafetyCoT 基础上引入 RLVR 后，安全得分提升至 $0.92$。虽然有用性略有下降，但整体仍维持在较高水平。这一变化反映出强化学习阶段对规则一致性的硬约束作用：当生成违反规则时，奖励函数将整体收益压制为零，从而对不安全行为形成更强惩罚。相比监督微调阶段的静态学习，RLVR 能够在分布外场景下进一步强化模型的边界稳定性。

综合来看，多模安全评估结果表明：仅依赖推理蒸馏无法稳定提升安全性，而将规则嵌入推理结构并结合强化验证机制，能够在安全与有用性之间形成更优的平衡点。

\subsection{越狱攻击鲁棒性分析}

\begin{table}[t]
\centering
\caption{越狱攻击安全评估结果。ASR 表示 Attack Success Rate，
Harm 表示有害程度平均得分（0--5，越低越好）。}
\label{tab:jailbreak}
\begin{tabular}{lcc}
\toprule
Model & ASR $\downarrow$ & Harm $\downarrow$ \\
\midrule
Base Model & 0.72 & 4.3 \\
Distilled CoT & 0.31 & 2.1 \\
SafetyCoT & 0.28 & 1.8 \\
SafetyCoT + RLVR & 0.12 & 1.1 \\
\bottomrule
\end{tabular}
\end{table}

在越狱攻击评估中，不同训练策略之间的性能差异更加显著。如表~\ref{tab:jailbreak} 所示，Base Model 的攻击成功率（ASR）高达 $0.72$，表明在模板化诱导与角色扮演类攻击下，模型容易被引导生成违规内容，同时 Harm 指标较高，说明生成内容具备较强的实际危害性。

Distilled CoT 显著降低了 ASR 至 $0.31$，说明显式推理链在一定程度上增强了模型对攻击意图的识别能力。然而，SafetyCoT 的 ASR 进一步下降至 $0.28$，且 Harm 值持续降低，说明规则约束下的安全反思训练能够提升模型在复杂语境中的边界判定稳定性。

值得注意的是，引入 RLVR 后，ASR 显著下降至 $0.12$，Harm 指标亦降至最低水平。这一结果说明强化学习阶段通过规则一致性奖励对越狱行为进行了显式惩罚，使模型在面对攻击性提示时形成更稳定的拒绝策略。从机制上看，RLVR 不仅优化了最终输出，还通过策略梯度更新强化了模型在内部推理阶段对风险信号的敏感性，从而提升了分布外攻击场景下的鲁棒性。

因此，越狱评估结果验证了本文方法在对抗性场景下的泛化能力，表明规则验证强化机制对于提升模型在复杂攻击模板下的稳定性具有关键作用。

\subsection{Agentic RAG 事实可信验证}

\begin{table*}[t]
\centering
\caption{Agentic RAG 事实可信评估结果。Fact Acc 表示与 Ground Truth 一致率，Safety 表示风险响应合规率，Latency 为平均响应时延。}
\label{tab:rag_results}
\begin{tabular}{lcccc}
\toprule
Model & Tool Trigger Acc $\uparrow$ & Fact Acc $\uparrow$ & Safety $\uparrow$ & Latency (ms) \\
\midrule
Vanilla RAG &  & 80.3  &  & 120 \\
Agentic RAG &  & 86.7 &  & 400 \\
\bottomrule
\end{tabular}
\end{table*}

在事实可信评估中，Vanilla RAG 能够在一定程度上提升事实覆盖度，但其工具调用行为缺乏风险约束机制，可能在检索到敏感或不完整信息时产生潜在不安全输出。相比之下，Agentic RAG 在 Fact Acc 上提升至 $86.7\%$，同时风险响应合规率同步提高，说明将安全反思能力迁移至工具调用与证据融合流程后，模型能够在事实生成前进行更为谨慎的决策。

从行为层面分析，Agentic RAG 不再将检索视为固定流程步骤，而是作为决策行为的一部分。当模型在内部推理阶段判断知识不确定时，能够主动触发工具调用，并结合规则边界对检索结果进行二次筛选。这种“安全反思—工具调用—证据融合”的闭环机制，有效缓解了检索增强带来的安全退化问题。

需要指出的是，引入多步决策与工具调用不可避免地增加了平均响应时延（Latency）。从实验结果来看，Agentic RAG 的时延显著高于 Vanilla RAG。这一现象反映了安全性、事实可靠性与实时性之间的工程权衡。在实际部署中，可根据业务场景对工具调用频率与触发阈值进行策略优化。